% RQ2 Background: Governance Capture

\subsection{Governance Capture in DAOs}

Governance capture occurs when actors gain disproportionate influence over collective decisions, steering outcomes toward private benefit rather than common good. In DAOs, capture typically manifests through:

\begin{itemize}
    \item \textbf{Token concentration}: Accumulating dominant token positions
    \item \textbf{Delegation capture}: Becoming a major delegate without proportional stake
    \item \textbf{Proposal manipulation}: Crafting proposals that benefit insiders
    \item \textbf{Vote timing}: Strategic voting at moments of low participation
\end{itemize}

\subsubsection{Empirical Evidence}

\citet{fritsch2022analyzing} documented extreme voting power concentration in MakerDAO, finding that a single proposal was decided by essentially one address. Similar patterns appear across protocols:

\begin{itemize}
    \item Uniswap: Top 10 addresses hold $>$40\% of voting power
    \item Compound: Venture capital firms control substantial delegated power
    \item Aave: Protocol treasury itself is a dominant voter
\end{itemize}

\subsection{Theoretical Foundations}

\subsubsection{Plutocracy vs. Democracy}

Token-weighted voting explicitly implements plutocracy: voting power proportional to wealth. Defenders argue this aligns incentives (stakeholders vote responsibly). Critics note it enables capture and discourages small holder participation.

The mechanism design challenge is creating systems that:
\begin{enumerate}
    \item Preserve stakeholder alignment (skin in the game)
    \item Limit maximum individual influence (anti-capture)
    \item Resist sybil attacks (can't trivially circumvent by splitting)
    \item Maintain operational efficiency (proposals can still pass)
\end{enumerate}

\subsubsection{The Sybil Problem}

Any per-address mitigation faces sybil attacks: actors splitting holdings across multiple addresses to circumvent limits. Perfect sybil resistance is impossible without identity systems, which conflict with pseudonymity norms.

Practical defenses include:
\begin{itemize}
    \item Making sybil attacks costly (gas fees, capital inefficiency)
    \item Using sublinear scaling (quadratic) rather than hard caps
    \item Temporal requirements (velocity penalties) that increase coordination cost
\end{itemize}

\subsection{Existing Mitigation Mechanisms}

\subsubsection{Quadratic Voting}

Introduced by \citet{weyl2018quadratic}, quadratic voting makes the marginal cost of additional votes increase linearly. Applied to DAOs, this typically means $\sqrt{\text{tokens}}$ voting power.

Gitcoin Grants uses quadratic funding (related mechanism) for public goods allocation. However, pure quadratic systems remain vulnerable to sybil attacks without identity verification.

\subsubsection{Vote Caps}

Some protocols implement maximum voting power per address. Optimism's Token House caps delegate power. MolochDAO variants use one-member-one-vote for certain decisions.

\subsubsection{Conviction Voting}

\citet{emmett2019conviction} proposed conviction voting where voting power accumulates over time. This naturally penalizes recent token acquisition and rewards long-term holders.

\subsubsection{Timelock and Vesting}

Protocols may require tokens to be locked or vested before gaining voting power. This is effectively a velocity penalty: recently acquired tokens have zero weight.

\subsection{Prior Simulation Work}

Limited prior work has systematically compared capture mitigation mechanisms. Theoretical analyses exist for quadratic voting \citep{lalley2018quadratic} but controlled empirical comparison is lacking. Our simulation framework addresses this gap.
